Über dem Alphabet $\Sigma=\{\texttt{0},\texttt{1}\}$ ist die Sprache
\[
L
=
\{
w\in\Sigma^*
\mid
|w|_{\texttt{0}}
-
|w|_{\texttt{1}}
\equiv
0
\mod 2
\}
\]
gegben.
Ist sie regulär?
Wenn ja, bestimmen Sie einen regulären Ausdruck dafür.

\begin{loesung}
Der endliche Automat 
\begin{center}
\begin{tikzpicture}[>=latex,thick]
\coordinate (q0) at (0,0);
\coordinate (q1) at (2,0);
\node at (q0) {$q_0$};
\node at (q1) {$q_1$};
\draw (q0) circle[radius=0.35];
\draw (q0) circle[radius=0.30];
\draw (q1) circle[radius=0.35];
\draw[->,shorten >= 0.35cm] ($(q0)+(-2,0)$) -- (q0);
\draw[->,shorten >= 0.35cm,shorten <= 0.35cm] (q0) to[out=30,in=150] (q1);
\draw[->,shorten >= 0.35cm,shorten <= 0.35cm] (q1) to[out=-150,in=-30] (q0);
\node at ($0.5*(q0)+0.5*(q1)+(0,0.3)$) [above] {\texttt{0},\texttt{1}};
\node at ($0.5*(q0)+0.5*(q1)+(0,-0.3)$) [below] {\texttt{0},\texttt{1}};
\end{tikzpicture}
\end{center}
akzeptiert die Sprache $L$, was beweist, dass $L$ regulär ist.

Ein regulärer Ausdruck kann nach dem Standardverfahren konstruiert werden.
Dazu wird der Automat erst in den folgenden VNEA standardisiert:
\begin{equation}
\raisebox{-1.3cm}{%
\begin{tikzpicture}[>=latex,thick]
\coordinate (s) at (-2,0);
\coordinate (q0) at (0,0);
\coordinate (q1) at (2,0);
\coordinate (e) at (0,-2);
\node at (q0) {$q_0$};
\node at (q1) {$q_1$};
\node at (e) {$e$};
\node at (s) {$s$};
\draw (s) circle[radius=0.35];
\draw (q0) circle[radius=0.35];
\draw (q1) circle[radius=0.35];
\draw (e) circle[radius=0.35];
\draw (e) circle[radius=0.30];
\draw[->,shorten >= 0.35cm] ($(s)+(-2,0)$) -- (s);
\draw[->,shorten >= 0.35cm,shorten <= 0.35cm] (s) -- (q0);
\draw[->,shorten >= 0.35cm,shorten <= 0.35cm] (q0) -- (e);
\draw[->,shorten >= 0.35cm,shorten <= 0.35cm] (q0) to[out=30,in=150] (q1);
\draw[->,shorten >= 0.35cm,shorten <= 0.35cm] (q1) to[out=-150,in=-30] (q0);
\node at ($0.5*(q0)+0.5*(q1)+(0,0.3)$) [above] {\texttt{0},\texttt{1}};
\node at ($0.5*(q0)+0.5*(q1)+(0,-0.3)$) [below] {\texttt{0},\texttt{1}};
\node at ($0.5*(s)+0.5*(q0)$) [above] {$\varepsilon$};
\node at ($0.5*(e)+0.5*(q0)$) [left] {$\varepsilon$};
\end{tikzpicture}}
\label{30000099:dea}
\end{equation}
Jetzt müssen die Zwischenzustände in der Reihenfolge $q_1,q_0$ entfernt werden,
es ergibt sich
\begin{center}
\definecolor{darkred}{rgb}{0.8,0,0}
\begin{tikzpicture}[>=latex,thick]
\begin{scope}
\coordinate (s) at (-2,0);
\coordinate (q0) at (0,0);
\coordinate (e) at (0,-2);
\node at (q0) {$q_0$};
\node at (e) {$e$};
\node at (s) {$s$};
\draw (s) circle[radius=0.35];
\draw (q0) circle[radius=0.35];
\draw (e) circle[radius=0.35];
\draw (e) circle[radius=0.30];
\draw[->,shorten >= 0.35cm] ($(s)+(-2,0)$) -- (s);
\draw[->,shorten >= 0.35cm,shorten <= 0.35cm] (s) -- (q0);
\draw[->,shorten >= 0.35cm,shorten <= 0.35cm] (q0) -- (e);
\draw[->,shorten >= 0.35cm,shorten <= 0.35cm]
	(q0) to[out=30,in=-30,distance=1.5cm] (q0);
\node at ($0.5*(s)+0.5*(q0)$) [above] {$\varepsilon$};
\node at ($0.5*(e)+0.5*(q0)$) [left] {$\varepsilon$};
\node at ($(q0)+(1.0,0.0)$) [right] {$\texttt{(0|1)(0|1)}$};
\end{scope}
\draw[->] (3.7,-1) -- +(1,0);
\begin{scope}[xshift=9cm]
\coordinate (s) at (-2,0);
\coordinate (e) at (0,-2);
\node at (e) {$e$};
\node at (s) {$s$};
\draw (s) circle[radius=0.35];
\draw (e) circle[radius=0.35];
\draw (e) circle[radius=0.30];
\draw[->,shorten >= 0.35cm] ($(s)+(-2,0)$) -- (s);
\draw[->,shorten >= 0.35cm,shorten <= 0.35cm] (s) -- (e);
\node at ($0.5*(s)+0.5*(e)$) [above right] {$\texttt{((0|1)(0|1))*}$};
\end{scope}
\end{tikzpicture}
\end{center}
Der gesuchte reguläre Ausdruck ist $r= \texttt{((0|1)(0|1))*}$.
Er kann kompakter auch $\texttt{(..)*}$ geschrieben werden.

Man kann auch direkt auf einen regulären Ausdruck für $L$ kommen.
Dazu beachtet man, dass die Bedingung bedeutet, dass das Wort $w\in L$
gerade Länge haben muss.
Dies ist auch, was der Automat \eqref{30000099:dea} ausdrückt.
Die Wörter gerader Länge werden aber auch vom regulären Ausdruck
$r=\texttt{(..)*}$ akzeptiert.
\end{loesung}

\begin{bewertung}
Endlicher Automat ({\bf E}) 1 Punkt,
Standardisierung ({\bf S}) 1 Punkt,
Entfernung von Zwischenzuständen ({\bf Z}) 2 Punkte,
finaler regulärer regulärer Ausdruck ({\bf R}) 1 Punkt,
Schlussfolgerung regulär ({\bf R}) 1 Punkt.
Bei Teillösungen: Idee, dass die Wortlänge gerade sein muss ({\bf G})
1--2 Punkte.
\end{bewertung}
